\chapter{Grafy planarne, ściany i wzór Eulera}
\label{ch:tg-planarne}

Dotychczas krawędzie grafu były abstrakcyjnymi odcinkami. Od tej chwili
istotne będzie również to, \emph{jak} graf umieszczamy w powierzchni.
Na przykład narysowanie dwóch krawędzi przecinających się na kartce
nie oznacza, że ich punkt przecięcia jest wierzchołkiem grafu.
Właśnie to rozróżnienie prowadzi do planarności, ścian, dualności
i do związku z charakterystyką Eulera powierzchni.

\section{Rysunek grafu a zanurzenie topologiczne}
\label{sec:tg-osadzanie-planarne}

\begin{definicja}[Zanurzenie grafu w powierzchnię]
Niech $G$ będzie skończonym grafem, dopuszczającym również pętle
i krawędzie równoległe, a $S$ powierzchnią topologiczną.
\emph{Zanurzeniem topologicznym grafu w $S$} (ang. \emph{embedding})
nazywamy ciągłe odwzorowanie różnowartościowe $\iota:|G|\to S$.
Z powodu zwartości $|G|$ i własności Hausdorffa przestrzeni $S$
jest ono homeomorfizmem na swój obraz. Graf nazywamy \emph{planarnym},
jeżeli istnieje takie $\iota$ dla $S=\mathbb R^2$.
W przypadku $S=\mathbb R^2$ konkretny graf wraz z wybranym
zanurzeniem nazywamy \emph{grafem płaskim}. Wnętrza obrazów różnych
krawędzi są rozłączne i nie zawierają obrazów wierzchołków;
dwie krawędzie mogą mieć wspólny punkt wyłącznie wtedy, gdy
mają wspólny koniec.
\end{definicja}

\begin{uwaga}[Płaszczyzna i sfera]
Graf skończony jest planarny wtedy i tylko wtedy, gdy można go
zanurzyć w sferze $S^2$. W jedną stronę stosujemy dołączenie
punktu w nieskończoności $\mathbb R^2\subset S^2$.
W drugą wybieramy punkt $p\in S^2\setminus\iota(|G|)$
(taki punkt istnieje, ponieważ skończona suma zanurzonych łuków
i punktów ma puste wnętrze w powierzchni)
i identyfikujemy $S^2\setminus\{p\}$ z $\mathbb R^2$.
Na sferze żadna ściana nie jest wyróżniona jako zewnętrzna.
\end{uwaga}

\begin{definicja}[Ściana]
Niech $\iota:|G|\hookrightarrow S^2$ będzie zanurzeniem grafu
skończonego. \emph{Ścianami} są składowe spójne dopełnienia
$S^2\setminus\iota(|G|)$. Ich liczbę oznaczamy przez $f(G,\iota)$.
Dla grafu płaskiego jedna ze ścian jest nieograniczona.
\end{definicja}

\begin{przyklad}[Ten sam graf, różne rysunki]
Graf $K_4$ można narysować z przecięciami na rysunku czworokąta
z obiema przekątnymi. Nie dowodzi to nieplanarności: wystarczy
umieścić trzy wierzchołki na brzegu trójkąta, a czwarty w jego wnętrzu.
Otrzymujemy wtedy cztery ściany na sferze, wszystkie ograniczone
trójkątami.
\end{przyklad}

\begin{center}
\begin{tikzpicture}[line cap=round,line join=round,scale=.95]
  \begin{scope}[shift={(-3.3,0)}]
    \coordinate (a) at (-1.15,-.9); \coordinate (b) at (1.15,-.9);
    \coordinate (c) at (0,1.30); \coordinate (d) at (0,-.15);
    \fill[tgBlue!8] (a)--(b)--(d)--cycle;
    \fill[tgTeal!12] (a)--(c)--(d)--cycle;
    \fill[tgGold!17] (b)--(c)--(d)--cycle;
    \draw[tgEdge] (a)--(b)--(c)--cycle (a)--(d)--(b) (c)--(d);
    \foreach \pt/\name in {a/a,b/b,c/c,d/d}
      \node[tgBlueVertex] at (\pt) {$\name$};
    \node[text=tgNavy,font=\small] at (0,-1.58) {$K_4$ bez przecięć};
  \end{scope}
  \begin{scope}[shift={(3.3,0)}]
    \coordinate (a) at (-1.10,-.95);\coordinate (b) at (1.10,-.95);
    \coordinate (c) at (1.10,.95);\coordinate (d) at (-1.10,.95);
    \draw[tgEdge] (a)--(b)--(c)--(d)--cycle;
    \draw[tgAccent] (a)--(c) (b)--(d);
    \foreach \pt/\name in {a/a,b/b,c/c,d/d}
      \node[tgRedVertex] at (\pt) {$\name$};
    \node[text=tgNavy,font=\small] at (0,-1.58) {inny rysunek $K_4$};
  \end{scope}
\end{tikzpicture}

{\footnotesize Czerwone przekątne po prawej przecinają się, lecz ten punkt
nie jest dodatkowym wierzchołkiem. Rysunek po lewej dowodzi planarności $K_4$.}
\end{center}

\section{Wzór Eulera i regularne otoczenia grafów}
\label{sec:tg-planarny-euler}

\begin{twierdzenie}[Euler dla grafów zanurzonych w sferze]
\label{thm:tg-euler-plane}
Niech $G$ będzie skończonym grafem o $n\ge1$ wierzchołkach,
$m$ krawędziach i $c$ składowych spójnych, zanurzonym w $S^2$.
Jeżeli $f$ oznacza liczbę ścian tego zanurzenia, to
\[
 n-m+f=1+c.
\]
W szczególności dla grafu spójnego zachodzi $n-m+f=2$.
Wynik jest prawdziwy także dla pętli i krawędzi równoległych.
\end{twierdzenie}

\begin{proof}
Korzystamy ze standardowej własności regularnych otoczeń
skończonych grafów osadzonych w powierzchni: wokół
wierzchołków można wybrać rozłączne dyski, a wzdłuż
krawędzi wąskie prostokątne pasy, z rozłącznymi wnętrzami,
przyklejone do odpowiednich łuków brzegów dysków.
Dla pętli oba łuki przyklejenia leżą na tym samym dysku
i są rozłączne. W przypadku osadzeń topologicznych
istnienie takiego otoczenia jest konsekwencją
twierdzenia Jordana--Schoenfliesa (lokalnego prostowania
skończonej rodziny łuków o wspólnym końcu).
Oznaczmy przez $N$ sumę dysków i pasów.
Jest to zwarta powierzchnia z brzegiem o $c$ składowych,
każda z nich jest podpowierzchnią sfery, a więc ma rodzaj zero.

Niech $b$ oznacza łączną liczbę okręgów brzegowych $N$.
Z jednej strony $\chi(N)=n-m$: zaczynamy od $n$ dysków,
każdy dołączany pas jest prostokątem o charakterystyce
Eulera $1$, przyklejanym wzdłuż dwóch rozłącznych
odcinków o łącznej charakterystyce $2$. Z addytywności
charakterystyki Eulera otrzymujemy zmianę $1-2=-1$.
Z drugiej strony każda zwarta spójna podpowierzchnia
sfery z $b_i$ okręgami brzegowymi jest sferą z usuniętymi
wnętrzami $b_i$ rozłącznych dysków (twierdzenie Jordana
dla okręgów brzegowych). Jej charakterystyka wynosi
$2-b_i$. Zsumowanie po składowych daje
\[
n-m=2c-b.
\]

Oznaczmy przez $M$ domknięcie dopełnienia $N$ w sferze.
Jego składowe odpowiadają ścianom grafu: w dyskach
wierzchołkowych dopełnienie grafu składa się z sektorów
między kolejnymi incydencjami, a w pasach z dwóch
połówek wzdłuż krawędzi. Sektory i połówki tworzą
kołnierze przy okręgach brzegowych $N$; dla izolowanego
wierzchołka jest to przebity dysk. Dołączenie tych
kołnierzy do składowych $M$ nie zmienia ich liczby.
W szczególności $M$ ma $f$ składowych, łącznie $b$
okręgów brzegowych i $\chi(M)=2f-b$.
Ponieważ $S^2=N\cup M$ i $N\cap M=\partial N$
jest sumą rozłączną okręgów o charakterystyce zero,
addytywność daje $2=(n-m)+(2f-b)$.
Podstawiając $b=2c-n+m$, otrzymujemy
$2=2(n-m+f-c)$, czyli tezę.
\end{proof}

\begin{uwaga}[Interpretacja przez dodawanie krawędzi]
\label{rem:tg-euler-edge-add}
W ustalonym osadzeniu dodanie krawędzi nie zmienia
liczby wierzchołków. Jeśli łączy ona różne składowe,
to $m$ rośnie o jeden, $c$ maleje o jeden, więc
ze wzoru Eulera $f$ się nie zmienia. Jeżeli jej końce
należą do tej samej składowej (w tym wypadek pętli),
to $c$ się nie zmienia i $f$ rośnie o jeden.
Wnętrze nowej krawędzi leży w jednej dotychczasowej
ścianie, więc tylko ta ściana może zostać podzielona.
Zatem w drugim przypadku dzieli się ona dokładnie na
dwie ściany. To uzasadnia krok indukcyjny często
stosowany w kombinatorycznym dowodzie wzoru Eulera.
\end{uwaga}


\begin{uwaga}[Co oznacza brzeg ściany]
Brzeg ściany nie musi być cyklem prostym. W drzewie
cała sfera poza grafem jest jedną ścianą, a przejście
wokół jej brzegu używa każdej krawędzi dwa razy.
W ogólności śledzimy \emph{obchody ścianowe}: zamknięte
spacery krawędziowe, w których liczymy każdą stronę
krawędzi osobno. Suma ich długości, liczonych z krotnościami,
wynosi $2m$, także gdy występują mosty i pętle.
Dla grafów niespójnych pojedyncza ściana może mieć kilka
obchodów brzegowych, a izolowany wierzchołek wnosi
obchód długości zero. Obchody odpowiadają okręgom
brzegowym regularnego otoczenia, niekoniecznie
różnym składowym topologicznym brzegu ściany.
\end{uwaga}

\begin{przyklad}[Drzewo, cykl i kwadrat z przekątną]
Każde drzewo o $n$ wierzchołkach ma $m=n-1$ krawędzi.
Euler daje $f=1$. Cykl $C_n$ ma $n$ krawędzi i dwie ściany:
wnętrze i zewnętrze. Kwadrat z jedną przekątną ma
$n=4$, $m=5$, więc $f=2-4+5=3$.
Dwie wewnętrzne ściany są trójkątne, trzecia jest zewnętrzna.
\end{przyklad}

\begin{center}
\begin{tikzpicture}[scale=.95,line cap=round,line join=round]
  \coordinate (a) at (-2,0);\coordinate (b) at (2,0);
  \coordinate (c) at (2,2.0);\coordinate (d) at (-2,2.0);
  \fill[tgBlue!11] (a)--(b)--(c)--cycle;
  \fill[tgTeal!11] (a)--(c)--(d)--cycle;
  \draw[tgEdge] (a)--(b)--(c)--(d)--cycle;
  \draw[tgAccent] (a)--(c);
  \foreach \pt in {a,b,c,d} \node[tgBlueVertex] at (\pt) {};
  \node[text=tgBlue] at (.9,.55) {$F_1$};
  \node[text=tgTeal] at (-.9,1.45) {$F_2$};
  \node[text=tgNavy] at (0,-.65) {$F_3$ -- ściana zewnętrzna};
\end{tikzpicture}
\end{center}

\section{Oszacowania liczby krawędzi i nieplanarność}
\label{sec:tg-planarne-granice}

\begin{lemat}[Stopnie ścianowe]
Niech $G$ będzie spójnym, prostym grafem płaskim,
$n\ge3$, o $m$ krawędziach i $f$ ścianach.
Długość obchodu brzegowego każdej ściany, liczona
z krotnościami skierowanych wystąpień krawędzi, jest
co najmniej trzy, a suma długości wszystkich obchodów
to $2m$. Jeśli $G$ jest ponadto dwudzielny, każdy
obchód jest parzystej długości, a więc ma długość co najmniej cztery.
\end{lemat}

\begin{proof}
Każda krawędź ma dwie strony, które należą do ścian
(niekoniecznie różnych). Idąc po obwodzie każdej ściany,
przechodzimy każdą stronę dokładnie raz, skąd suma długości
wynosi $2m$. Obchód długości jeden wymagałby pętli, co jest
sprzeczne z prostotą. Obchód długości dwa mógłby powstać
z dwóch krawędzi równoległych albo z pojedynczej krawędzi
przechodzonej dwukrotnie. Jeśli zaś obchód składa się z dwukrotnego przejścia
jednej krawędzi, krawędź ta jest mostem, a obydwa jej końce
mają stopień jeden w całym grafie: obchód ściany
przechodzi od danej incydencji do kolejnej w porządku
cyklicznym wokół wierzchołka. Natychmiastowy powrót
tą samą krawędzią wymaga, aby była ona jedyną
incydencją przy tym wierzchołku.
W przypadku grafu spójnego oznaczałoby to, że graf jest
pojedynczą krawędzią $K_2$; wyklucza to $n\ge3$.
Zatem każda długość jest co najmniej trzy.

W grafie dwudzielnym dowolny zamknięty spacer przechodzi
naprzemiennie pomiędzy obiema klasami wierzchołków,
ma więc długość parzystą. Wobec poprzedniego oszacowania
długość każdej ściany jest nie mniejsza niż cztery.
\end{proof}

\begin{twierdzenie}[Ograniczenia planarne]
\label{thm:tg-planarne-bound}
Jeżeli $G$ jest spójnym grafem prostym, planarnym,
o $n\ge3$ wierzchołkach i $m$ krawędziach, to
$m\le3n-6$. Jeżeli ponadto $G$ jest dwudzielny,
to $m\le2n-4$.
\end{twierdzenie}

\begin{proof}
Ustalmy zanurzenie płaskie. Poprzedni lemat daje
$3f\le2m$. Ze wzoru Eulera $f=2-n+m$,
więc $3(2-n+m)\le2m$, co po uporządkowaniu
daje $m\le3n-6$. W przypadku grafu dwudzielnego
mamy mocniejsze $4f\le2m$, a stąd
$4(2-n+m)\le2m$ oraz $m\le2n-4$.
\end{proof}

\begin{przyklad}[Dwie klasyczne przeszkody]
Graf $K_5$ ma pięć wierzchołków i dziesięć krawędzi,
lecz $10>3\cdot5-6=9$, więc nie jest planarny.
Graf $K_{3,3}$ ma sześć wierzchołków i dziewięć krawędzi,
jest dwudzielny, a $9>2\cdot6-4=8$, więc również
nie jest planarny. Samo ograniczenie $m\le3n-6$
nie wystarczyłoby do wykrycia nieplanarności $K_{3,3}$.
\end{przyklad}

\begin{center}
\begin{tikzpicture}[scale=.82,line cap=round,>=Stealth]
 \begin{scope}[shift={(-3.3,0)}]
  \foreach \i in {0,...,4} {
    \coordinate (p\i) at ({90-72*\i}:1.2);
  }
  \foreach \i in {0,...,4} {
    \foreach \j in {0,...,4} {
      \ifnum\i<\j \draw[tgEdge] (p\i)--(p\j);\fi
    }
  }
  \foreach \i in {0,...,4} \node[tgBlueVertex] at (p\i) {};
  \node[text=tgNavy,font=\small] at (0,-1.80) {$K_5$: $10>9$};
 \end{scope}
 \begin{scope}[shift={(3.3,0)}]
  \foreach \i in {1,2,3} {
    \coordinate (u\i) at (-1.25,{2-\i});
    \coordinate (v\i) at (1.25,{2-\i});
  }
  \foreach \i in {1,2,3} \foreach \j in {1,2,3}
      \draw[tgEdge] (u\i)--(v\j);
  \foreach \i in {1,2,3} {
    \node[tgBlueVertex] at (u\i) {};
    \node[tgRedVertex] at (v\i) {};
  }
  \node[text=tgNavy,font=\small] at (0,-1.80) {$K_{3,3}$: $9>8$};
 \end{scope}
\end{tikzpicture}

{\footnotesize Rysunki mają przecięcia niebędące wierzchołkami.
Nieplanarność tych grafów wynika z dowodów liczbowych, nie z wyglądu rysunków.}
\end{center}

\section{Kolorowanie ścian i graf dwudzielny w planie}
\label{sec:tg-planarne-color}

\begin{twierdzenie}[Kryterium dwukolorowalności ścian]
\label{thm:tg-face-two-color}
Niech $G$ będzie skończonym grafem o niepustym zbiorze
wierzchołków, osadzonym w $S^2$, z dopuszczalnymi pętlami
i krawędziami równoległymi (nie zakładamy spójności).
Ściany można pokolorować dwoma kolorami tak, aby
po obu stronach każdej krawędzi występowały różne
kolory, wtedy i tylko wtedy, gdy każdy wierzchołek
ma parzysty stopień. Pętlę liczymy w stopniu dwukrotnie.
\end{twierdzenie}

\begin{proof}
Jeżeli takie kolorowanie istnieje, rozważmy mały okrąg
w dysku otaczającym dowolny nieizolowany wierzchołek $v$.
Przechodzimy kolejno przez sektory między incydencjami,
zmieniając kolor dokładnie raz przy każdej incydencji.
Po pełnym obiegu wracamy do wyjściowego sektora, więc
liczba zmian, równa $d_G(v)$, musi być parzysta.
Wierzchołki izolowane mają stopień zero.

Odwrotnie, załóżmy parzystość wszystkich stopni.
Zbiór krawędzi grafu można rozłożyć na rozłączne
krawędziowo cykle proste (wliczając pętle i cykle
z dwóch krawędzi równoległych). Istotnie, w każdej
niepustej składowej krawędziowej istnieje zamknięty
szlak Eulera. Fragment pomiędzy pierwszym powtórzeniem
wierzchołka daje cykl prosty. Usuwamy jego krawędzie,
co zmniejsza stopnie należących do niego wierzchołków
o dwa. Powtarzamy aż do wyczerpania krawędzi.
Jeżeli graf nie ma krawędzi, rodzina cykli jest pusta.

Niech $\Gamma_1,\ldots,\Gamma_k$ będą obrazami
otrzymanych cykli. Każdy jest krzywą Jordana.
Wybierzmy dla każdego $i$ jedną ze składowych $U_i$
dopełnienia $S^2\setminus\Gamma_i$. Dla punktu
$x$ poza grafem określamy kolor
\[
\kappa(x)=\sum_{i=1}^k\chi_{U_i}(x)\pmod 2.
\]
Na każdej ścianie wszystkie funkcje $\chi_{U_i}$
są stałe: ściana jest spójna i nie przecina żadnej
krzywej $\Gamma_i$. Dlatego kolor zależy wyłącznie
od ściany. Wybierzmy punkt we wnętrzu krawędzi $e$.
Dokładnie jedna krzywa $\Gamma_i$ zawiera tę krawędź.
W dostatecznie małym otoczeniu punktu pozostałe krzywe
są nieobecne, a dwa lokalne boki $\Gamma_i$ należą
do różnych składowych jej dopełnienia, na mocy
twierdzenia Jordana--Schoenfliesa. Przejście przez $e$
zmienia więc dokładnie jeden składnik sumy, czyli kolor.
Otrzymujemy wymagane kolorowanie.
\end{proof}

\begin{wniosek}[Most i ściany]
\label{cor:tg-bridge-face}
W skończonym grafie osadzonym w sferze krawędź
jest mostem wtedy i tylko wtedy, gdy po obu jej
stronach występuje ta sama ściana.
W szczególności graf z parzystymi stopniami
wszystkich wierzchołków nie zawiera mostów.
\end{wniosek}
\begin{proof}
Po usunięciu wnętrza krawędzi mogą się połączyć
wyłącznie ściany sąsiadujące z tą krawędzią:
w regularnym otoczeniu usunięcie pasa otwiera
przejście między jego dwiema stronami.
Jeżeli są to różne ściany, ich liczba maleje o jeden,
a jeżeli ta sama, pozostaje bez zmiany. Z uwagi
\ref{rem:tg-euler-edge-add} drugi przypadek zachodzi
dokładnie wtedy, gdy dodanie krawędzi do grafu bez niej
scala dwie składowe, czyli gdy krawędź jest mostem.
Ostatnia teza wynika z twierdzenia
\ref{thm:tg-face-two-color}.
\end{proof}

\begin{uwaga}
Warunek parzystych stopni wierzchołków nie jest równoważny
dwudzielności grafu: cykl trójkątny ma wszystkie stopnie
parzyste, lecz nie jest dwudzielny. Stwierdzenie dotyczy
kolorowania \emph{ścian}, a nie wierzchołków.
\end{uwaga}

\begin{przyklad}[Dwukolorowanie ścian bez dwudzielności grafu]
Niech $G$ powstaje z dwóch trójkątów przez
utożsamienie dokładnie jednego ich wierzchołka $v$.
Wówczas $d_G(v)=4$, a pozostałe cztery wierzchołki
mają stopień $2$. Graf jest więc eulerowski,
ale nie jest dwudzielny, gdyż zawiera trójkąty.
Jego standardowe osadzenie na sferze ma trzy ściany:
dwie trójkątne oraz zewnętrzną.
Obu ścianom trójkątnym możemy przypisać kolor $A$,
a ścianie zewnętrznej kolor $B$.
Każda krawędź oddziela wówczas ściany różnych kolorów,
co ilustruje poprzednie stwierdzenie.
\end{przyklad}

\begin{center}
\begin{tikzpicture}[scale=.96,line cap=round,line join=round]
  \coordinate (o) at (0,0);
  \coordinate (a) at (-1.65,.95);
  \coordinate (b) at (-1.65,-.95);
  \coordinate (c) at (1.65,.95);
  \coordinate (d) at (1.65,-.95);
  \fill[tgTeal!20] (o)--(a)--(b)--cycle;
  \fill[tgTeal!20] (o)--(c)--(d)--cycle;
  \draw[tgEdge] (o)--(a)--(b)--cycle
                 (o)--(c)--(d)--cycle;
  \foreach \pt/\n in {o/v,a/a,b/b,c/c,d/d}
    \node[tgBlueVertex] at (\pt) {$\n$};
  \node[text=tgNavy,font=\footnotesize] at (-1.07,0) {$A$};
  \node[text=tgNavy,font=\footnotesize] at (1.07,0) {$A$};
  \node[text=tgNavy,font=\footnotesize] at (0,1.42)
    {ściana zewnętrzna: $B$};
\end{tikzpicture}

{\footnotesize Dwa trójkąty mają wspólny wierzchołek $v$.
Ściany trójkątne mają kolor $A$, a zewnętrzna kolor $B$.}
\end{center}


% ============================================================
% Iteracja 4, rozdział 10
% ============================================================
